% ===========================================================================
\section{The ring equation: limbs, quotient, and product witness}\label{sec:arith}
% ===========================================================================

\subsection{Trace layout}\label{sec:layout}

The trace has $2^{\nu}$ rows, \emph{one signature per row}. Fix the
\emph{cell width} $W = 32$ (the native commitment-cell degree bound) and
the \emph{limb count} $L = n/W = 16$. Each degree-$<n$ ring element is
stored as $L$ \emph{limb cells}, each a degree-$<W$ \texttt{arbitrary\_poly}
cell holding $W$ consecutive integer coefficients. A polynomial
$P(X) = \sum_{j=0}^{n-1} P[j] X^j$ is recovered from its limbs
$\{\col{P}_m\}_{m<L}$, $\col{P}_m(X) = \sum_{p<W} P[mW+p]X^p$, by the
\emph{limb reconstruction}
\begin{equation}\label{eq:recon}
  P(X) \;=\; \sum_{m=0}^{L-1} X^{W m}\, \col{P}_m(X).
\end{equation}
Keeping cells at the native width $W=32$ (rather than one giant
degree-$n$ cell) lets the \emph{same} cells feed both the ring equation
(via \cref{eq:recon}) and the norm zerocheck (\cref{sec:norm}, which reads
the coefficients off the limb cells directly) --- no second
representation, no consistency bridge. Crucially, the trace cell degree
$W$ is \emph{decoupled} from the ring degree $n$: the scalars $X^{Wm}$
used in \cref{eq:recon} live in a separate scalar ring whose degree bound
is independent of $W$.

\smallskip
\noindent The columns, all \texttt{arbitrary\_poly} unless noted, are:
\begin{center}
\begin{tabular}{@{}llll@{}}
\toprule
columns & role & kind & limbs \\
\midrule
$\col{h}_{0..L}$       & shared public key (identical on every row) & public  & $L$ \\
$\col{c}_{0..L}$       & per-signature hash point                   & public  & $L$ \\
$\col{s_1}_{0..L}$     & recomputed short-vector half (centered)    & witness & $L$ \\
$\col{s_2}_{0..L}$     & signature short-vector half                & witness & $L$ \\
$\col{u}_{0..L}$       & mod-$q$ / negacyclic quotient (Option B)   & witness & $L$ \\
$\col{s_3}_{0..L_3}$   & product witness $s_2 h$ (full, unreduced)  & witness & $L_3 = 2L$ \\
$\col{slack}$          & per-signature squared-norm slack ($\Z$ column) & witness & --- \\
\bottomrule
\end{tabular}
\end{center}
Public columns ($\col{h}$, $\col{c}$) are absorbed into the transcript;
$\col{h}$ is shared, so it is one public polynomial broadcast to every
row, while $\col{c}$ varies per row. With $5L + L_3 = 112$
\texttt{arbitrary\_poly} columns and one $\Z$ column.

\subsection{Option B: the ring equation over \texorpdfstring{$\Z[X]$}{Z[X]}}\label{sec:optionb}

The Zinc+ UAIR supports \emph{ideal-membership} constraints
$Q(\cdot)\in\ideal{g}$ over $\Z[X]$ (equality is the special case
$g=0$). The negacyclic ring is the quotient by
$g = X^n+1$, realized in the implementation as the rotation ideal
$\langle X^{W'} - a\rangle$ with $W' = n$ and $a = -1$, whose membership
test \texttt{RotationIdeal::remainder\_is\_zero} reduces a polynomial of
\emph{any} degree modulo $X^n+1$ in $O(n)$ by Horner folding (so degrees
up to $<2n-1$ are handled natively).

The reduction modulo $q$ in \cref{eq:ring} is made explicit by an integer
\emph{quotient} witness $u$ (``Option B''): \cref{eq:ring} holds in $\Rq$
iff there is $u\in\Rring$ with
$s_1 + s_2 h - c = q\,u$ in $\Rring$, i.e.
\begin{equation}\label{eq:ring-z}
  s_1 + s_2 h - c - q\,u \;\in\; \negacyc \quad\text{over } \Z[X].
\end{equation}
For the honest prover, $u$ is the centered quotient of the (negacyclically
reduced) value $s_1 + (s_2 h \bmod X^n{+}1) - c$ by $q$; its coefficients
fit comfortably in a machine word.

\subsection{The product witness: keeping the ring equation linear}\label{sec:s3}

Written directly, \cref{eq:ring-z} contains the product $s_2 h$ of two
\emph{trace} polynomials. The constraint degree counter (which assigns
degree $1$ to each trace cell and adds degrees on multiplication) reports
this constraint at \emph{degree~$2$}. A degree-$2$ effective constraint
forces the prover onto the slow \emph{combined-polynomial} ideal-check
lane; the fast \emph{MLE-first} lane (column-major, no per-row combined
polynomial) requires every non-zero-ideal constraint to be linear.

We therefore carry the product as a committed witness. Introduce
$\col{s_3}$ holding the \textbf{full, unreduced} product $s_3 := s_2 h$
over $\Z[X]$ (degree $<2n-1$, hence $L_3 = 2L = 32$ limbs), and split the
constraint into two:
\begin{align}
  \text{\textbf{ring equation}}\quad
    & s_1 + s_3 - c - q\,u \;\in\; \negacyc,
      \label{eq:ring-linear}\\[2pt]
  \text{\textbf{product witness}}\quad
    & s_3 - s_2 h \;=\; 0.
      \label{eq:product}
\end{align}
In \cref{eq:ring-linear} the term $\col{s_3}$ is a trace value (recovered
from its $L_3$ limbs by \cref{eq:recon}), \emph{not} a product, so the
residual is \emph{degree~$1$ --- linear}. \cref{eq:product} is the only
nonlinear constraint, and we assert it as a \textbf{zero ideal}
(\texttt{assert\_zero}). This is sound and complete precisely because
$\col{s_3}$ carries the \emph{full} product: $s_3 - s_2 h$ is then
\emph{identically} the zero polynomial over $\Z[X]$, not merely a multiple
of $X^n+1$. (Had $\col{s_3}$ stored the negacyclically reduced product, it
would differ from $s_2 h$ by $(X^n{+}1)\cdot(\text{quotient})$ and the
\texttt{assert\_zero} would be incomplete.)

\subsection{Why this is MLE-first eligible}\label{sec:mlefirst}

The Zinc+ degree machinery distinguishes \emph{two} degrees:
\begin{itemize}[leftmargin=1.5em]
  \item \texttt{count\_max\_degree} --- the maximum per-variable degree
  over \emph{all} constraints; drives the inner sumcheck, which runs at
  \texttt{count\_max\_degree}$+2$.
  \item \texttt{count\_effective\_max\_degree} --- the same maximum but
  \emph{excluding zero-ideal} constraints; drives the ideal-check lane
  selection (MLE-first requires it to be $\le 1$).
\end{itemize}
For our instance the ring equation \eqref{eq:ring-linear} is degree~$1$
(non-zero ideal) and the product \eqref{eq:product} is degree~$2$ (zero
ideal). Hence
\[
  \texttt{count\_effective\_max\_degree} = 1, \qquad
  \texttt{count\_max\_degree} = 2,
\]
so the instance is \textbf{MLE-first eligible}, while the inner sumcheck
runs at degree $2+2 = 4$.

\begin{remark}[Soundness of the zero-ideal split]\label{rem:soundness}
Excluding \eqref{eq:product} from the \emph{effective} degree does
\emph{not} drop it from the proof. Its ideal-check value is forced to the
only member of the zero ideal --- $0$ --- so the verifier's membership
check is satisfied vacuously; but the inner sumcheck (at degree~$4$)
folds the \emph{actual} expression $s_3 - s_2 h$ against the committed
$\col{s_2},\col{h},\col{s_3}$. A cheating prover with $s_3 \ne s_2 h$ on
some row produces a sumcheck sum that disagrees with the (zero) ideal-check
value and is rejected. This is the same pattern the SHA-256 UAIR uses for
its degree-$2$ \texttt{assert\_zero} selector pins; here the quadratic term
$s_2 h$ is a witness $\times$ public product, structurally identical to
SHA's $\col{s\_msg\_init}\cdot(\col{w\_W}-\col{pa\_m})$.
\end{remark}

\begin{remark}[A subtlety in the MLE-first lane]\label{rem:mlefirst-zero}
The MLE-first ideal check evaluates each constraint on the
\emph{point-evaluated} column MLEs. For a degree-$2$ zero-ideal term this
yields $s_3^\ast - s_2^\ast h^\ast$ (a product of MLE evaluations), which
is \emph{nonzero} at a random point even though the constraint vanishes on
the hypercube ($\mle{s_3 - s_2 h}$ is the MLE of the product, not the
product of MLEs). The implementation must therefore force the zero-ideal
slot to $0$ before transcript absorption; otherwise the ideal-check value
and the sumcheck disagree. This is enforced in
\texttt{IdealCheckProtocol::prove\_linear}.
\end{remark}
